\section{Supplementary Token-Score Diagnostics}
\label{app:token_scores}
We compare token scores as supplementary diagnostics at the final 1.7B and 4B checkpoints (Tables~\ref{tab:paired_selectors} and~\ref{tab:paired4b_selectors}). Each score selects the same fraction of positions. The tables report captured KL and hidden-gradient energy; these scores do not alter the recorded training updates.

\paragraph{Disagreement, uncertainty, and log-ratio variance.}
All scores use the conditional distributions in Equation~\ref{eq:loss}. We record KL, total variation $\frac12\sum_v|\tilde p_v-\tilde q_v|$, and student entropy $H=-\sum_v\tilde p_v\log\tilde p_v$. For $r_v=\log(\tilde p_v/\tilde q_v)$, the log-ratio variance is $V=\operatorname{Var}_{\tilde p}(r_v)$. This variance appears as the reverse-KL local utility in \citet{sadagopan2026utility} and as the Fisher-metric signal term in \citet{sheng2026ier}.

\paragraph{TIP-inspired entropy--KL combination.}
Following the Soft-OR construction of \citet{xu2026tip}, we compute $s=\bar H+\bar K-\bar H\bar K$, where bars denote batchwise min--max normalization. Our measurements apply this score to the recorded Top-16 conditional distributions. They evaluate the score under our checkpoint protocol; the paper's complete token-selection training procedure is outside this comparison.

\paragraph{Information-efficiency ratio.}
Following \citet{sheng2026ier}, we evaluate the conditional-support information-efficiency ratio (IER)
\begin{equation}
 b=\frac{\sum_v(1-\tilde p_v)r_v}{\sum_v(1-\tilde p_v)},\qquad
 \operatorname{IER}=\frac{V}{\sum_v(1-\tilde p_v)(r_v-b)^2-V}.
\end{equation}
We rank valid positive-denominator cases. IER concerns a sampled next-token gradient estimator; our training directly evaluates the selected-support loss. Thus, the table measures a hypothetical estimator's reliability under the cached conditional distributions.

\paragraph{Representation-dependent diagnostics.}
We additionally define retention $\rho=\|U^\top g_h\|^2/\|g_h\|^2$ and the combined score $s\rho$ to inspect the measured subspace. These are geometric diagnostics used in this study. At both scales, disagreement and log-ratio variance capture more full-hidden gradient energy than entropy; retention alone captures little total energy, while the combination shifts selection toward projected energy. A training intervention would be needed to establish benefits from using these rankings during optimization.

\begin{table*}[!htbp]
\centering
\caption{Signal captured by the highest-scoring decile at the final paired-study checkpoint. Each cell reports FULL. Columns give captured KL, full hidden-gradient energy, and $U$-projected energy, all in percent. Scores are observational diagnostics.}
\label{tab:paired_selectors}
\small
\setlength{\tabcolsep}{3pt}
\begin{tabular}{lrrr}
\toprule
Selection score & KL & Full energy & $U$ energy \\
\midrule
KL & 77.06 & 84.69 & 84.36 \\
Total variation & 74.37 & 78.81 & 77.83 \\
Entropy & 37.37 & 38.29 & 29.04 \\
Entropy--KL mixture & 40.68 & 42.71 & 34.76 \\
Log-ratio variance & 71.91 & 86.78 & 86.33 \\
Estimator reliability & 59.17 & 80.25 & 76.53 \\
$U$ retention & 11.70 & 7.56 & 21.17 \\
Mixture $\times$ retention & 37.65 & 38.62 & 63.92 \\
\bottomrule
\end{tabular}
\end{table*}

\begin{table*}[!htbp]
\centering
\caption{4B. Signal captured by the highest-scoring decile at the final paired-study checkpoint. Each cell reports FULL. Columns give captured KL, full hidden-gradient energy, and $U$-projected energy, all in percent. Scores are observational diagnostics.}
\label{tab:paired4b_selectors}
\small
\setlength{\tabcolsep}{3pt}
\begin{tabular}{lrrr}
\toprule
Selection score & KL & Full energy & $U$ energy \\
\midrule
KL & 74.79 & 81.87 & 79.06 \\
Total variation & 72.49 & 79.71 & 77.50 \\
Entropy & 39.22 & 40.04 & 31.85 \\
Entropy--KL mixture & 43.47 & 44.46 & 37.04 \\
Log-ratio variance & 70.11 & 83.34 & 80.63 \\
Estimator reliability & 59.99 & 76.14 & 70.79 \\
$U$ retention & 6.66 & 5.23 & 16.22 \\
Mixture $\times$ retention & 32.37 & 32.84 & 61.36 \\
\bottomrule
\end{tabular}
\end{table*}


\FloatBarrier
\section{Displacement and Basis Discovery}
\label{app:displacement}
Figure~\ref{fig:rq2_adaptation} uses the paired FULL trajectories, matching the main supervision and gradient analyses. Tables~\ref{tab:paired_adapter_geometry} and~\ref{tab:paired4b_adapter_geometry} provide their endpoint values. The separate original studies replicate the displacement/gradient coverage contrast and show strong overlap between nearby late-checkpoint bases (Table~\ref{tab:subspace}). For rank-$r$ orthonormal bases, overlap is $\|U_a^\top U_b\|_F^2/r$.




\begin{table}[t]
\centering
\caption{Geometry at the final FULL checkpoint. Own-$r$ rows give cumulative displacement energy captured by that checkpoint\textquotesingle s leading directions. Fixed-$U$ uses the step-10 basis. The last two rows measure subspace overlap. All entries are percentages.}
\label{tab:subspace}
\small
\setlength{\tabcolsep}{3pt}
\begin{tabular}{lrr}
\toprule
Metric (\%) & 1.7B & 4B \\
\midrule
Own-1 motion & 49.58 & 23.97 \\
Own-4 motion & 72.03 & 55.70 \\
Own-8 motion & 76.56 & 63.94 \\
Own-16 motion & 80.68 & 70.20 \\
Own-128 motion & 90.55 & 84.88 \\
Fixed-$U$ motion & 88.01 & 80.10 \\
Fixed-$U$ gradient & 16.79 & 15.18 \\
$U_{10},U_{106}$ overlap & 77.86 & 72.17 \\
$U_{100},U_{106}$ overlap & 99.46 & 98.72 \\
\bottomrule
\end{tabular}
\end{table}

\begin{table*}[!htbp]
\centering
\caption{1.7B paired FULL final displacement and gradient geometry. Own-$r$ coverage uses final displacement eigenvectors; fixed-$U$ uses the step-10 basis. Energy shares are percentages.}
\label{tab:paired_adapter_geometry}
\small
\setlength{\tabcolsep}{3pt}
\begin{tabular}{lr}
\toprule
Measurement & Value \\
\midrule
FULL final own-1 / own-128 & 53.02 / 90.36 \\
FULL final fixed-$U$ coverage & 87.46 \\
FULL final fixed-$U$ gradient energy & 16.71 \\
Hidden dimension / isotropic rank-128 share (\%) & 2048 / 6.25 \\
\bottomrule
\end{tabular}
\end{table*}

\begin{table*}[!htbp]
\centering
\caption{4B paired FULL final displacement and gradient geometry. Own-$r$ coverage uses final displacement eigenvectors; fixed-$U$ uses the step-10 basis. Energy shares are percentages.}
\label{tab:paired4b_adapter_geometry}
\small
\setlength{\tabcolsep}{3pt}
\begin{tabular}{lr}
\toprule
Measurement & Value \\
\midrule
FULL final own-1 / own-128 & 22.24 / 84.14 \\
FULL final fixed-$U$ coverage & 79.16 \\
FULL final fixed-$U$ gradient energy & 15.18 \\
Hidden dimension / isotropic rank-128 share (\%) & 2560 / 5.00 \\
\bottomrule
\end{tabular}
\end{table*}


\begin{table}[t]
\centering
\caption{First accepted stopping window. Checks occur every five updates, using three observations and starting no earlier than update 15. Overlaps and energy gap are percentages. Required minima: adjacent overlap 98\%, endpoints 97\%, lower-decile squared principal-angle cosine 90\%; maximum energy gap: 0.5 percentage points.}
\label{tab:stopping}
\small
\setlength{\tabcolsep}{3pt}
\begin{tabular}{lrr}
\toprule
Criterion & 1.7B & 4B \\
\midrule
Stop update & 40 & 50 \\
Adjacent min. (\%) & 98.70 & 98.43 \\
Endpoints (\%) & 97.96 & 97.40 \\
Lower decile (\%) & 96.21 & 95.49 \\
Energy gap (pp) & 0.13 & 0.14 \\
\bottomrule
\end{tabular}
\end{table}


\FloatBarrier
\subsection{Cumulative Displacement Magnitudes and Early-Basis Coverage}
\label{app:cumulative_magnitudes}
Table~\ref{tab:cumulative_motion_audit} reports every paired FULL motion observation from basis discovery onward. The fixed panels contain 5,203 positions for 1.7B and 6,481 for 4B. Both cumulative displacement magnitudes and early-basis coverage are retained, so the coverage percentages can be read alongside the absolute scale of the measured change.

\begin{table}[!htbp]
\centering
\caption{Cumulative displacement magnitude and fixed early-basis coverage at every recorded checkpoint from step 10 in paired FULL. RMS is the square root of the mean squared vector norm across fixed probe positions. Raw RMS uses unnormalized hidden states; unit RMS uses unit-normalized states. Coverage is measured on the unit-normalized displacement relative to initialization, using the fixed step-10 rank-128 basis. All entries describe cumulative changes.}
\label{tab:cumulative_motion_audit}
\small\setlength{\tabcolsep}{8pt}
\begin{tabular}{rrrrrrr}
\toprule
 & \multicolumn{3}{c}{1.7B} & \multicolumn{3}{c}{4B} \\
\cmidrule(lr){2-4}\cmidrule(lr){5-7}
Step & Raw RMS & Unit RMS & $U$ coverage (\%) & Raw RMS & Unit RMS & $U$ coverage (\%) \\
\midrule
10 & 47.357 & 0.33499 & 85.66 & 57.784 & 0.36566 & 80.50 \\
15 & 56.045 & 0.40512 & 86.50 & 62.306 & 0.39057 & 80.92 \\
20 & 58.821 & 0.42436 & 86.57 & 63.665 & 0.40182 & 80.59 \\
25 & 62.165 & 0.44992 & 86.79 & 64.463 & 0.40821 & 80.29 \\
30 & 61.876 & 0.44841 & 86.48 & 65.053 & 0.41356 & 80.09 \\
35 & 60.366 & 0.43684 & 86.10 & 65.947 & 0.42071 & 79.80 \\
40 & 61.393 & 0.44512 & 86.32 & 66.744 & 0.42587 & 79.72 \\
45 & 62.098 & 0.45074 & 86.51 & 67.666 & 0.43196 & 79.61 \\
50 & 62.610 & 0.45513 & 86.64 & 67.905 & 0.43286 & 79.62 \\
55 & 63.720 & 0.46361 & 86.92 & 67.924 & 0.43262 & 79.59 \\
60 & 64.293 & 0.46756 & 87.00 & 68.456 & 0.43603 & 79.46 \\
65 & 64.854 & 0.47148 & 87.08 & 68.694 & 0.43763 & 79.39 \\
70 & 66.268 & 0.48192 & 87.29 & 68.937 & 0.43856 & 79.27 \\
75 & 65.890 & 0.47819 & 87.24 & 68.876 & 0.43779 & 79.21 \\
80 & 66.648 & 0.48370 & 87.30 & 69.016 & 0.43847 & 79.24 \\
85 & 67.046 & 0.48624 & 87.24 & 69.241 & 0.43977 & 79.18 \\
90 & 67.386 & 0.48835 & 87.27 & 69.404 & 0.44056 & 79.19 \\
95 & 67.485 & 0.48898 & 87.24 & 69.804 & 0.44260 & 79.08 \\
100 & 67.865 & 0.49170 & 87.28 & 69.852 & 0.44290 & 79.06 \\
105 & 68.632 & 0.49771 & 87.44 & 70.200 & 0.44534 & 79.14 \\
106 & 68.707 & 0.49835 & 87.46 & 70.269 & 0.44572 & 79.16 \\
\bottomrule
\end{tabular}
\end{table}

